% 第 16 讲 Tate 上同调群的消没（译自 src/lecture16.txt）
\section*{第 16 讲\quad Tate 上同调群的消没}

前情回顾：我们把（上同调的）局部类域论归约成了如下命题：设 $L/K$ 是非阿基米德局部域的有限 Galois 扩张，Galois 群为 $G$，则
\[
\big((L\otimes_K K^{\mathrm{unr}})^\times\big)_{tG}\simeq 0,
\]
即该复形无环。为证明这一消没，我们将证明一个（属于 Tate 的）关于 Tate 上同调消没的一般定理，它使上述问题变得更加易于处理。于是我们问：给定 $G$-模的复形 $X$，什么条件能保证 $X_{tG}$ 无环？这类结果的原型如下：

\textbf{定理 16.1.} 若 $G$ 是循环群，则 $X_{tG}$ 无环当且仅当
\[
\widehat H^0(G,X)=0=\widehat H^1(G,X).
\]

\textit{证明.} $X_{tG}$ 可由一个 $2$-周期复形来计算。注意任意两个奇偶性不同的值（例如相邻的两个值）就足够了。\hfill$\square$

本讲的主要结果是：

\textbf{定理 16.2.} 若 $G$ 是 $p$-群（即 $\#G=p^n$，其中 $p$ 为某个素数、$n\ge 0$），定理 16.1 仍然成立。

由此我们将推出下一个结果：

\textbf{定理 16.3.} 设对每个素数 $p$ 及 $G$ 的每个 $p$-Sylow 子群 $G_p\subseteq G$，都有 $\widehat H^0(G_p,X)=0=\widehat H^1(G_p,X)$。则 $X_{tG}$ 无环。

\textbf{注 16.4.} 一般而言，对任意的有限群 $G$，「在相邻两个次数上消没」并不充分。此外，实践中人们往往是对 $G$ 的每个子群（而不只是 $p$-Sylow 子群）验证相邻两个次数上 Tate 上同调的消没。

在接下来的几讲中，我们将由此推出局部类域论。我们先证明定理 16.2。以下始终设 $G$ 为 $p$-群。

\textbf{命题 16.5.} 设 $X$ 是 $\mathbb F_p[G]$-模的复形。若 $\widehat H^0(G,X)=0$，则 $X_{tG}$ 无环。

注意这里我们只需要一个次数上的消没！为此，我们先回忆下述事实。

\textbf{引理 16.6.} 设 $V$ 是非零的 $\mathbb F_p[G]$-模。则 $V^G\neq 0$。

\textit{证明.} 不妨设 $V$ 在 $\mathbb F_p$ 上有限维，因为 $V$ 显然是它的有限维 $G$-子模的并。于是 $\#V=p^r$，其中 $0<r<\infty$。$V$ 的每个 $G$-轨道的大小或者为 $1$、或者被 $p$ 整除（轨道大小必定整除 $\#G$，因为该轨道同构于 $G$ 对稳定化子的商）。由于 $\{0\}$ 是一个大小为 $1$ 的 $G$-轨道，必存在另一个大小为 $1$ 的轨道（因为所有 $G$-轨道的大小之和必须等于 $p^r$），即存在某个 $v\in V^G\setminus\{0\}$。\hfill$\square$

\textit{证明（命题 16.5）.} \textbf{第 1 步.} 先断言：对每个有限维 $\mathbb F_p$-表示 $V$，都有 $\widehat H^0(G,X\otimes V)=0$。这里 $G$ 通过「对角作用」作用，即在 $(X\otimes V)^i=X^i\otimes V$ 上由 $g\cdot(x\otimes v):=gx\otimes gv$ 给出。我们对 $\dim_{\mathbb F_p}V$ 作归纳。由前一引理，有 $\mathbb F_p[G]$-模的短正合列
\[
0\to V^G\to V\to W\to 0,
\]
其中 $\dim W<\dim V$。这给出
\[
\operatorname{hCoker}(V^G\otimes X\to V\otimes X)\simeq W\otimes X.
\]
由于由归纳假设 $\widehat H^0(G,W\otimes X)=0$，而
\[
\widehat H^0(G,V^G\otimes X)=\widehat H^0\Big(G,\bigoplus_{\dim V^G}X\Big)=\bigoplus_{\dim V^G}\widehat H^0(G,X)=0
\]
（由对 $X$ 的假设），Tate 上同调的长正合列便给出 $\widehat H^0(G,V\otimes X)=0$，如所欲证。

\textbf{第 2 步.} 现在证明负次数上的消没。考虑短正合列
\[
0\to V_1\to \mathbb F_p[G]\xrightarrow{\ \varepsilon\ }\mathbb F_p\to 0,
\]
其中 $V_1$ 定义为 $\varepsilon$ 的核，这与我们把 $\mathbb F_p$ 换成 $\mathbb Z$ 时所说的 $I_G$ 类似。设 $\underline X$ 表示 $X$ 配上平凡的 $G$-作用。则
\[
\underline X\otimes\mathbb F_p[G]\to X\otimes\mathbb F_p[G],\qquad x\otimes g\mapsto gx\otimes g
\]
（右端取对角作用）是 $G$-等变映射，并且是双射，从而是同构。事实上，
\[
h(x\otimes g)=x\otimes hg\mapsto hgx\otimes hg=h(gx\otimes g).
\]
习题集 7 第 1(e) 题已经证明 $(\underline X\otimes\mathbb F_p[G])_{tG}$ 无环，故 $(X\otimes\mathbb F_p[G])_{tG}$ 亦然。于是 Tate 上同调的长正合列给出
\[
\widehat H^{i-1}(G,X)\simeq \widehat H^i(G,X\otimes V_1).
\]
第 1 步已表明右端在 $i=0$ 时消没，因此 $\widehat H^{-1}(G,X)=0$。迭代即得 $\widehat H^i(G,X)=0$ 对一切 $i\le 0$ 成立。

\textbf{第 3 步.} 为证明正次数上的消没，注意我们有正合列
\[
0\to \mathbb F_p\xrightarrow{\ 1\mapsto\sum_{g\in G}g\ }\mathbb F_p[G]\to V_2\to 0,
\]
其中 $V_2$ 如前面一样定义为余核。同样的论证给出
\[
\widehat H^i(G,X\otimes V_2)\simeq \widehat H^{i+1}(G,X),
\]
于是第 1 步再次表明 $\widehat H^i(G,X)=0$ 对一切 $i\ge 0$ 成立。\hfill$\square$

\textit{证明（定理 16.2）.} 定义 $X/p:=\operatorname{hCoker}(X\xrightarrow{p}X)$；注意这与把每一项都对 $p$ 取商并不相同。注意，作为 $\mathbb Z[G]$-模的复形，$X/p$ 拟同构于一个 $\mathbb F_p[G]$-模的复形。由于 $X$ 只在拟同构的意义下有定义，不妨设它作为 $\mathbb Z[G]$-模的复形是投射的（特别地，平坦的）。于是我们有拟同构
\[
X/p=X\otimes_{\mathbb Z[G]}\operatorname{hCoker}(\mathbb Z[G]\xrightarrow{p}\mathbb Z[G])\simeq X\otimes_{\mathbb Z[G]}\mathbb F_p[G],
\]
其中右端逐项计算 $X$ 模 $p$ 的商。于是有长正合列
\[
\cdots\to \widehat H^i(G,X)\xrightarrow{p}\widehat H^i(G,X)\to \widehat H^i(G,X/p)\to \widehat H^{i+1}(G,X)\xrightarrow{p}\widehat H^{i+1}(G,X)\to\cdots,
\]
令 $i=0$，由假设即得 $\widehat H^0(G,X/p)=0$。于是由命题 16.5，$(X/p)_{tG}$ 无环。由此 $\widehat H^i(G,X/p)=0$ 对一切 $i$ 成立，因此乘以 $p$ 在 $\widehat H^i(G,X)$ 上是同构，对每个 $i$ 皆然。但正如习题集 7 第 2(c) 题所示，乘以 $\#G$ 在 $\widehat H^i(G,X)$ 上是零映射。由于 $G$ 是 $p$-群，这只有在 $\widehat H^i(G,X)=0$ 对一切 $i$ 成立时才可能。\hfill$\square$

\textit{证明（定理 16.3）.} 如上所述，乘以 $\#G$ 在 $\widehat H^i(G,X)$ 上是零映射，故 $\widehat H^i(G,X)$ 是 $\#G$-挠的。因此只需证明 $\widehat H^i(G,X)[p]=0$ 对一切 $p$ 成立（即 $\widehat H^i(G,X)$ 的 $p$-挠部分消没）。

回忆对 $G$ 的每个子群 $H$，有分别称为限制与余限制的映射 $X_{tG}\to X_{tH}$ 与 $X_{tH}\to X_{tG}$，其复合作为 $X_{tG}$ 的自同态同伦于乘以指数 $[G:H]$。% [原文疑误：原印作 restriction 与 inflation（膨胀），但复合为乘以指数 $[G:H]$ 的正是「限制与余限制」（corestriction），且膨胀映射只对正规子群有定义，故此处径改为余限制]

对 $p$-Sylow 子群 $H=G_p$ 应用上述结论并取上同调，我们得到映射
\[
\widehat H^i(G,X)[p]\subset \widehat H^i(G,X)\to \widehat H^i(G_p,X)\to \widehat H^i(G,X),
\]
其复合是乘以 $[G:G_p]$，而后者按定义与 $p$ 互素。因此它在限制到 $\widehat H^i(G,X)[p]$ 上是同构，特别地 $\widehat H^i(G,X)[p]\to\widehat H^i(G_p,X)$ 是单射。但由定理 16.2，$\widehat H^i(G_p,X)=0$ 对一切 $i$ 成立，于是得到所欲结论。\hfill$\square$
